\begin{table}[htbp]
\centering
\small
\begin{tabular}{lrrrrrrr}
\toprule
 & 0 & 1 & 2 & 3 & 4 & 5 & 7 \\
\midrule
0 & 1{,}00 & 0{,}13 & 0{,}05 & -0{,}04 & 0{,}02 & 0{,}03 & -0{,}10 \\
1 & 0{,}13 & 1{,}00 & 0{,}53 & 0{,}23 & 0{,}64 & 0{,}45 & -0{,}27 \\
2 & 0{,}05 & 0{,}53 & 1{,}00 & 0{,}57 & 0{,}38 & 0{,}99 & 0{,}11 \\
3 & -0{,}04 & 0{,}23 & 0{,}57 & 1{,}00 & 0{,}15 & 0{,}61 & 0{,}43 \\
4 & 0{,}02 & 0{,}64 & 0{,}38 & 0{,}15 & 1{,}00 & 0{,}35 & -0{,}16 \\
5 & 0{,}03 & 0{,}45 & 0{,}99 & 0{,}61 & 0{,}35 & 1{,}00 & 0{,}16 \\
7 & -0{,}10 & -0{,}27 & 0{,}11 & 0{,}43 & -0{,}16 & 0{,}16 & 1{,}00 \\
\bottomrule
\end{tabular}
\caption{A normált forgatókönyvek járási allokációinak Spearman‑rangkorrelációja.}
\label{tab:spearman}
\forras{saját számítás az \texttt{analysis/} pipeline kimeneteiből}
\end{table}
